% HTML-specific semantic preamble for LaTeXML.  The print preamble is deliberately
% not loaded here: KOMA page layout, tcolorbox skins, headers, and minitocs add a
% large amount of TeX state without contributing useful structure to the web book.
\usepackage{fontspec}
\setmainfont{Noto Serif CJK KR}
\setsansfont{Noto Sans CJK KR}
\setmonofont{Noto Sans Mono CJK KR}

\usepackage{amsmath,amssymb,amsthm}
\usepackage{mathrsfs,mathtools,stmaryrd,wasysym}
\usepackage[usenames,svgnames,dvipsnames]{xcolor}
\usepackage{enumerate,multirow,graphicx}
\usepackage{hyperref}
\usepackage[nameinlink]{cleveref}
\usepackage{tikz-cd}
\tikzcdset{surjective head/.style={two heads}}

\input{tex/macros}

\newcommand{\prototype}[1]{\par\noindent\textbf{대표적인 예:} #1\par}
\newcommand{\problemhead}{생각해 볼 만한 조금 더 어려운 문제}
\newcommand{\listhack}{}
\newcommand{\parttoc}{}
\newcommand{\subtitle}[1]{\gdef\napkinsubtitle{#1}}
\newcommand{\napkinsubtitle}{}
\newcommand{\twodigits}[1]{\ifnum#1<10 0\fi\the#1}
\newcommand{\chili}{\includegraphics[width=1.5em]{media/chili.png}}
\newcommand{\onechili}{\chili}
\newcommand{\twochili}{\chili\chili}
\newcommand{\threechili}{\chili\chili\chili}
\newcommand{\odif}{\mathop{}\!\mathrm{d}}
\newcommand{\pdv}[2]{\frac{\partial #1}{\partial #2}}
\newcommand{\todo}[2][]{\par\noindent\textbf{작성 예정:} #2\par}
\newcommand{\missingfigure}[1]{%
  \begin{quote}\textbf{그림 작성 예정:} #1\end{quote}}

\newenvironment{moral}{\begin{quote}\bfseries}{\end{quote}}
\newenvironment{hint}{\begin{quote}\textbf{힌트.} }{\end{quote}}
\newenvironment{sol}{\begin{quote}\textbf{풀이.} }{\end{quote}}
\newenvironment{tcolorbox}[1][]{\begin{quote}}{\end{quote}}

\newtheorem{theorem}{정리}[section]
\newtheorem{lemma}[theorem]{보조정리}
\newtheorem{proposition}[theorem]{명제}
\newtheorem{corollary}[theorem]{따름정리}
\theoremstyle{definition}
\newtheorem{example}[theorem]{예제}
\newtheorem{definition}[theorem]{정의}
\newtheorem{claim}[theorem]{주장}
\newtheorem{fact}[theorem]{사실}
\newtheorem{abuse}[theorem]{표기의 남용}
\newtheorem{ques}[theorem]{질문}
\newtheorem{exercise}[theorem]{연습문제}
\newtheorem{step}[theorem]{단계}
\theoremstyle{remark}
\newtheorem{remark}[theorem]{비고}
\newtheorem*{remark*}{비고}
\newtheorem{problem}{문제}[chapter]
\newtheorem{sproblem}[problem]{문제}
\newtheorem{dproblem}[problem]{문제}
\renewcommand{\theproblem}{\thechapter\Alph{problem}}
\renewcommand{\thesproblem}{\theproblem$^{\star}$}
\renewcommand{\thedproblem}{\theproblem$^{\dagger}$}

\newcommand{\napkinkoreancrefname}[2]{%
  \crefname{#1}{#2}{#2}\Crefname{#1}{#2}{#2}}
\newcommand{\napkinkoreannames}{%
  \renewcommand{\contentsname}{목차}%
  \renewcommand{\partname}{부}%
  \renewcommand{\chaptername}{장}%
  \renewcommand{\figurename}{그림}%
  \renewcommand{\tablename}{표}%
  \renewcommand{\bibname}{참고 문헌}%
  \renewcommand{\proofname}{증명}%
  \napkinkoreancrefname{part}{부}%
  \napkinkoreancrefname{chapter}{장}%
  \napkinkoreancrefname{section}{절}%
  \napkinkoreancrefname{subsection}{소절}%
  \napkinkoreancrefname{figure}{그림}%
  \napkinkoreancrefname{table}{표}%
  \napkinkoreancrefname{equation}{식}%
  \napkinkoreancrefname{theorem}{정리}%
  \napkinkoreancrefname{lemma}{보조정리}%
  \napkinkoreancrefname{proposition}{명제}%
  \napkinkoreancrefname{corollary}{따름정리}%
  \napkinkoreancrefname{definition}{정의}%
  \napkinkoreancrefname{example}{예제}%
  \napkinkoreancrefname{remark}{비고}%
  \napkinkoreancrefname{ques}{질문}%
  \napkinkoreancrefname{exercise}{연습문제}%
  \napkinkoreancrefname{claim}{주장}%
  \napkinkoreancrefname{fact}{사실}%
  \napkinkoreancrefname{abuse}{표기의 남용}%
  \napkinkoreancrefname{problem}{문제}%
  \napkinkoreancrefname{sproblem}{문제}%
  \napkinkoreancrefname{dproblem}{문제}%
  \napkinkoreancrefname{step}{단계}%
  \napkinkoreancrefname{item}{항목}}
